% Рамки для формул. Подключается в преамбуле description_ru.tex и
% report_2026-11/report.tex после xcolor.
%   paperbox{заголовок} --- формула статьи как напечатано (серая);
%   resbox[заголовок]   --- наш результат (синяя);
%   cmpbox[заголовок]   --- сравнение со статьёй, вердикт (оранжевая);
%   keybox              --- ключевая формула без заголовка (синяя рамка).
\usepackage[most]{tcolorbox}
\tcbset{kdbase/.style={breakable,enhanced,arc=1mm,boxrule=0.5pt,
  left=2mm,right=2mm,top=1mm,bottom=1mm,before skip=6pt,after skip=6pt,
  fonttitle=\small\bfseries,toptitle=0.5mm,bottomtitle=0.5mm}}
\newtcolorbox{paperbox}[1]{kdbase,colback=gray!4,colframe=gray!55!black,
  colbacktitle=gray!18,coltitle=black,title={#1}}
\newtcolorbox{resbox}[1][Результат]{kdbase,colback=blue!3,
  colframe=blue!45!black,colbacktitle=blue!12,coltitle=black,title={#1}}
\newtcolorbox{cmpbox}[1][Сравнение со статьёй]{kdbase,colback=orange!4,
  colframe=orange!65!black,colbacktitle=orange!15,coltitle=black,title={#1}}
\newtcolorbox{keybox}{kdbase,colback=blue!2,colframe=blue!40!black}
